% SPDX-License-Identifier: MIT
% Copyright (c) 2026 Evgeny Baulin
% Optimization Methods, Seminar 06: theory handout (Russian).
% Generated by checks/06. Second-Order Methods Part Two.ipynb from web/06/js/content-*.js and content/students.json. Do not edit by hand.
% Compile with: latexmk -xelatex -interaction=nonstopmode -halt-on-error Theory_ru.tex
\documentclass[11pt,a4paper]{article}

\usepackage{fontspec}
\setmainfont{cmunrm.otf}[BoldFont=cmunbx.otf, ItalicFont=cmunti.otf, BoldItalicFont=cmunbi.otf]
\setsansfont{cmunss.otf}[BoldFont=cmunsx.otf, ItalicFont=cmunsi.otf, BoldItalicFont=cmunso.otf]
\setmonofont{cmuntt.otf}[BoldFont=cmuntb.otf, ItalicFont=cmunit.otf, BoldItalicFont=cmuntx.otf]
\usepackage{polyglossia}
\setdefaultlanguage{russian}
\setotherlanguage{english}
\newfontfamily\cyrillicfont{cmunrm.otf}[Ligatures=TeX, BoldFont=cmunbx.otf, ItalicFont=cmunti.otf, BoldItalicFont=cmunbi.otf]
\newfontfamily\cyrillicfontsf{cmunss.otf}[Ligatures=TeX, BoldFont=cmunsx.otf, ItalicFont=cmunsi.otf, BoldItalicFont=cmunso.otf]
\newfontfamily\cyrillicfonttt{cmuntt.otf}[BoldFont=cmuntb.otf, ItalicFont=cmunit.otf, BoldItalicFont=cmuntx.otf]

\usepackage[a4paper,left=2.3cm,right=2.3cm,top=2.4cm,bottom=2.4cm,headheight=14pt]{geometry}
\usepackage{amsmath,amssymb}
\usepackage{graphicx}
\usepackage{array,booktabs,longtable}
\usepackage{enumitem}
\usepackage{xcolor}
\usepackage[most]{tcolorbox}
\usepackage{titlesec}
\usepackage{fancyhdr}
\usepackage{needspace}
\usepackage[expansion=false]{microtype}
\usepackage[colorlinks=true,linkcolor=hseblue,urlcolor=hseblue,citecolor=hseblue]{hyperref}

\definecolor{hseblue}{RGB}{16,45,105}
\definecolor{hseaccent}{RGB}{198,40,40}
\definecolor{hsegray}{RGB}{108,114,126}
\definecolor{hselight}{RGB}{240,243,248}
\definecolor{hsewarm}{RGB}{253,244,232}

\hypersetup{pdftitle={Методы оптимизации. Семинар 6. Конспект теории}, pdfauthor={Евгений Баулин}}

\titleformat{\section}{\color{hseblue}\Large\bfseries}{\thesection}{0.6em}{}
\titleformat{\subsection}{\color{hseblue}\large\bfseries}{\thesubsection}{0.6em}{}
\titlespacing*{\section}{0pt}{16pt plus 4pt minus 2pt}{8pt plus 2pt}
\titlespacing*{\subsection}{0pt}{12pt plus 3pt minus 2pt}{6pt plus 2pt}
\pretocmd{\section}{\Needspace{7\baselineskip}}{}{\errmessage{cannot patch section}}
\pretocmd{\subsection}{\Needspace{5\baselineskip}}{}{\errmessage{cannot patch subsection}}

\tcbset{handoutbox/.style={enhanced, breakable, sharp corners, boxrule=0pt, left=7pt, right=7pt, top=5pt, bottom=5pt, before skip=8pt plus 2pt, after skip=8pt plus 2pt}}
\newtcolorbox{examplebox}{handoutbox, colback=hselight, colframe=hselight, before upper={\textcolor{hseblue}{\bfseries Пример.}\ }}
\newtcolorbox{whybox}{handoutbox, colback=white, colframe=hseblue, leftrule=2.5pt, top=3pt, bottom=3pt, before upper={\textcolor{hseaccent}{\bfseries Почему.}\ }}
\newtcolorbox{answerbox}{handoutbox, colback=hsewarm, colframe=hsewarm, before upper={\textcolor{hseaccent}{\bfseries Ответ.}\ }}
\newtcolorbox{conventionbox}{handoutbox, colback=hselight, colframe=hseblue, leftrule=2.5pt, before upper={\textcolor{hseblue}{\bfseries Обозначения.}\par\smallskip}}

\newcommand{\handoutfigure}[3]{\par\medskip\noindent\begin{minipage}{\linewidth}\centering
  \includegraphics[width=#1\linewidth]{figures/#2.pdf}\par\smallskip
  \begin{minipage}{0.94\linewidth}\small\textcolor{hseblue}{\bfseries Рисунок.}\ #3\end{minipage}\end{minipage}\par\medskip}
\newcommand{\problemhead}[3]{\par\bigskip\noindent{\color{hseblue}\large\bfseries #1.\ #2}\par\nopagebreak
  \noindent{\small\itshape\color{hsegray}#3}\par\nopagebreak\smallskip}
\newtcolorbox{definitionbox}[1]{handoutbox, colback=white, colframe=hseblue, leftrule=2.5pt, before upper={\textcolor{hseblue}{\bfseries #1}\ }}
\newtcolorbox{theorembox}[1]{handoutbox, colback=hselight, colframe=hseblue, leftrule=2.5pt, before upper={\textcolor{hseblue}{\bfseries #1}\ }}
\newtcolorbox{proofbox}{handoutbox, colback=white, colframe=hsegray, leftrule=1.2pt, top=3pt, bottom=3pt, before upper={\textcolor{hsegray}{\bfseries\itshape Доказательство.}\ }}
\newtcolorbox{notebox}[1]{handoutbox, colback=hselight, colframe=hselight, before upper={\textcolor{hseblue}{\bfseries #1}\par\smallskip}}
\newtcolorbox{trapbox}{handoutbox, colback=hsewarm, colframe=hseaccent, leftrule=2.5pt, before upper={\textcolor{hseaccent}{\bfseries Ловушка.}\ }}
\newtcolorbox{selfcheckbox}{handoutbox, colback=white, colframe=hseaccent, leftrule=2.5pt, top=3pt, bottom=3pt, before upper={\textcolor{hseaccent}{\bfseries Проверьте себя.}\ }}
\newcommand{\handoutdisplay}[1]{\sbox0{$\displaystyle #1$}\ifdim\wd0>\linewidth
  \par{\centering\binoppenalty=10000\relpenalty=0\relax$\displaystyle #1$\par}\else\[#1\]\fi}
\newcommand{\selfcheckanswer}{\par\smallskip\textcolor{hsegray}{\bfseries Ответ.}\ }

\setlist[itemize]{itemsep=3pt,topsep=4pt}
\setlist[enumerate]{itemsep=3pt,topsep=4pt}
\setlength{\parindent}{0pt}
\setlength{\parskip}{5pt plus 1pt}
\setlength{\emergencystretch}{2em}
\setcounter{tocdepth}{2}

\pagestyle{fancy}
\fancyhf{}
\fancyhead[L]{\small\color{hsegray}Методы оптимизации}
\fancyhead[R]{\small\color{hsegray}Семинар 6}
\fancyfoot[L]{\small\color{hsegray}Евгений Баулин}
\fancyfoot[C]{\small\thepage}
\fancyfoot[R]{\small\color{hsegray}Лицензия MIT}
\renewcommand{\headrulewidth}{0pt}
\renewcommand{\footrulewidth}{0pt}
\fancypagestyle{plain}{\pagestyle{fancy}}

\begin{document}

\begin{center}
  {\large\color{hsegray}Методы оптимизации\par}
  \vspace{4pt}
  {\large\bfseries\color{hseaccent}Семинар 6\par}
  \vspace{8pt}
  {\LARGE\bfseries\color{hseblue}Методы второго порядка, часть вторая: Гаусс–Ньютон, Левенберг–Марквардт, доверительная область, L-BFGS\par}
  \vspace{8pt}
  {\Large Конспект теории\par}
  \vspace{12pt}
  НИУ ВШЭ, факультет компьютерных наук\par
  2026-2027 учебный год\par
  \vspace{8pt}
  Лектор курса: Андрей Игнатов\par
  Евгений Баулин\par
  Telegram: \href{https://t.me/tarakan_tuc}{@tarakan\_tuc}\par
\end{center}

\noindent{\color{hseblue}\rule{\linewidth}{0.6pt}}

\tableofcontents

\clearpage


Методы второго порядка, продолжение. На семинаре 5 были шаг Ньютона и его защита. Конспект следует лекции и компьютерному семинару Андрея Игнатова о методах Гаусса–Ньютона и Левенберга–Марквардта для наименьших квадратов, о доверительной области как другой защите и об L-BFGS как квазиньютоновском методе с ограниченной памятью.

\begin{conventionbox}
$r:\mathbb R^n\to\mathbb R^m$ — вектор невязок с якобианом $J$. $B$ — текущая оценка гессиана, $G$ — обратной к нему; $s_k=x_{k+1}-x_k$, $y_k=g_{k+1}-g_k$. Радиус доверия — $\Delta$; квадратичная модель — $m(p)=f+g^\top p+\tfrac{1}{2}p^\top Bp$, а $\rho=(\text{фактическая})/(\text{предсказанная})$.
\end{conventionbox}

% =====================================================================
\section{Метод Гаусса–Ньютона}\label{sec:gn}

Для $f(x)=\tfrac{1}{2}\|r(x)\|^2$ с $r:\mathbb R^n\to\mathbb R^m$ и якобианом $J$ имеем $\nabla f=J^\top r$ и $\nabla^2 f=J^\top J+\sum_i r_i\nabla^2 r_i$. Метод Гаусса–Ньютона оставляет только $J^\top J$ и решает $J^\top Jp=-J^\top r$.

\begin{whybox}
Шаг Гаусса–Ньютона — это шаг Ньютона для линеаризованной невязки $r+Jp$. Отброшенное слагаемое — то, которое требовало бы вторых производных $r$; оно мало, когда невязки в решении малы.
\end{whybox}

\begin{theorembox}{Теорема: Сходимость метода Гаусса–Ньютона.}
Если $J(x^*)$ имеет полный столбцовый ранг и $r(x^*)=0$, существует окрестность $x^*$, в которой Гаусс–Ньютон сходится к $x^*$ квадратично. Если же $\|r(x^*)\|$ отделено от нуля, сходимость в лучшем случае линейная.
\end{theorembox}

\begin{trapbox}
$J^\top J$ положительно полуопределена, но не определена: если $J$ не полного ранга, нормальные уравнения становятся вырожденными. Левенберг–Марквардт решает это сдвигом; QR-разложение $J$ — более точная альтернатива, никогда не строящая $J^\top J$.
\end{trapbox}

\problemhead{A1}{Гаусс–Ньютон на подгонке экспоненты}{У доски}
\textbf{Условие.} Подгоните $f_\theta(t)=a\,e^{bt}$ к четырём точкам $(t_i,y_i)=\begin{pmatrix} 0 & 2 \\ 1 & 3{,}3 \\ 2 & 5{,}4 \\ 3 & 9 \end{pmatrix}$ с $\theta=(a,b)$. Из $\theta_0=(1;\ 1)$ сделайте два шага Гаусса–Ньютона. Укажите $\theta_1$, $\theta_2$ и $\tfrac{1}{2}\|r(\theta_2)\|^2$.

\textbf{Решение.}
\begingroup\sloppy
\begin{enumerate}[label=\textbf{\arabic*.}, leftmargin=*]
  \item В точке $\theta_0$ модельные значения $a_0e^{b_0t_i}=(1;\ 2{,}7183;\ 7{,}3891;\ 20{,}0855)$, значит, $r_0=(-1;\ -0{,}5817;\ 1{,}9891;\ 11{,}0855)$. У якобиана строки $(e^{b_0t_i},a_0t_ie^{b_0t_i})$: $J_0=\begin{pmatrix} 1 & 0 \\ 2{,}7183 & 2{,}7183 \\ 7{,}3891 & 14{,}7781 \\ 20{,}0855 & 60{,}2566 \end{pmatrix}$.
  \item Нормальные уравнения $J_0^\top J_0\,p=-J_0^\top r_0$ дают $J_0^\top J_0=\begin{pmatrix} 466{,}416 & 1326{,}8717 \\ 1326{,}8717 & 3856{,}6408 \end{pmatrix}$ и $J_0^\top r_0=(234{,}7749;\ 695{,}7901)$. Решая по правилу Крамера, получаем $p_0=(0{,}4654;\ -0{,}3405)$, значит, $\theta_1=\theta_0+p_0=(1{,}4654;\ 0{,}6595)$.
  \item Повторяем в $\theta_1$: модель $(1{,}4654;\ 2{,}8337;\ 5{,}4798;\ 10{,}5967)$, невязки $(-0{,}5346;\ -0{,}4663;\ 0{,}0798;\ 1{,}5967)$, $J_1=\begin{pmatrix} 1 & 0 \\ 1{,}9338 & 2{,}8337 \\ 3{,}7395 & 10{,}9595 \\ 7{,}2314 & 31{,}79 \end{pmatrix}$. Нормальные уравнения дают $p_1=(0{,}4557;\ -0{,}1548)$ и $\theta_2=(1{,}9211;\ 0{,}5047)$. Сумма квадратов падает с $64{,}0919$ до $1{,}5295$ и далее до $0{,}0542$: уже после второго шага мы возле минимума.
  \item Точное второе слагаемое гессиана равно $\sum_i r_i\nabla^2 r_i$: при малых невязках оно мало, и Гаусс–Ньютон ведёт себя как Ньютон. Точки лежат на экспоненциальной кривой, невязки быстро уменьшаются, а сходимость почти квадратичная.
\end{enumerate}
\endgroup
\begin{answerbox}
$\theta_1=(1{,}4654;\ 0{,}6595)$, $\theta_2=(1{,}9211;\ 0{,}5047)$, $\tfrac{1}{2}\|r(\theta_2)\|^2=0{,}0542$.
\end{answerbox}

% =====================================================================
\section{Метод Левенберга–Марквардта}\label{sec:lm}

\begin{definitionbox}{Определение.}
\textbf{Метод Левенберга–Марквардта} решает $(J^\top J+\lambda I)p=-J^\top r$ с $\lambda\ge0$, которое адаптируется по отношению $\rho=\frac{f(x)-f(x+p)}{m(0)-m(p)}$: при хорошем $\rho$ $\lambda$ уменьшается, при плохом — растёт.
\end{definitionbox}

Малое $\lambda$ — это Гаусс–Ньютон; большое — короткий градиентный шаг $-J^\top r/\lambda$. Метод интерполирует между ними без необходимости выбирать стратегию линейного поиска.

\problemhead{A2}{Левенберг–Марквардт спасает дикое начало}{Самостоятельно}
\textbf{Условие.} В задаче подгонки из A1 начните из далёкой точки $\theta_0=(10;\ 1{,}5)$ и сделайте три шага Левенберга–Марквардта с $\lambda_0=1$. Правило: $\rho=(f(\theta)-f(\theta+p))/(m(0)-m(p))$; если $\rho\ge0{,}25$, шаг принимается и $\lambda\leftarrow\lambda/10$; иначе шаг отвергается и $\lambda\leftarrow\lambda\cdot10$. Укажите $\theta$ после первого и после третьего принятого шага, а также число отказов.

\textbf{Решение.}
\begingroup\sloppy
\begin{enumerate}[label=\textbf{\arabic*.}, leftmargin=*]
  \item Из $\theta_0$ полный шаг Гаусса–Ньютона уходит слишком далеко; LM с $\lambda_0=1$ его укорачивает. Первые пробные шаги имеют $\rho<0{,}25$ и отвергаются; $\lambda$ каждый раз увеличивается в десять раз, пока сдвига не станет достаточно.
  \item После первого принятого шага $\theta=(0{,}8765;\ 1{,}474)$. Три принятых шага приводят $\theta$ к $(1{,}0785;\ 0{,}8463)$; суммарное число отказов — $0$.
  \item Каждое принятие делит $\lambda$ на десять, каждое отвержение — умножает на десять. Метод начинается как градиентный, по мере роста доверия к модели ведёт себя всё более как Гаусс–Ньютон и заканчивается очень близко к его пределу.
\end{enumerate}
\endgroup
\begin{answerbox}
После одного принятого шага $a=0{,}8765$; после трёх, $b=0{,}8463$; отказов — $0$.
\end{answerbox}

% =====================================================================
\section{Доверительная область}\label{sec:tr}

\begin{definitionbox}{Определение.}
\textbf{Шаг доверительной области} минимизирует квадратичную модель при $\|p\|\le\Delta$:
\handoutdisplay{
p=\arg\min_{\|p\|\le\Delta} f+g^\top p+\tfrac{1}{2}p^\top Bp.
}
 По ККТ минимум удовлетворяет $(B+\lambda I)p=-g$ при некотором $\lambda\ge0$ с условием $\lambda(\|p\|-\Delta)=0$.
\end{definitionbox}

\textbf{Dogleg} приближает решение ломаной $0\to p_C\to p_N$: берут самую дальнюю точку этой ломаной, лежащую в шаре. Это точно, когда $B\succ0$ и путь dogleg совпадает с кривой решений доверительной области.

\begin{itemize}
  \item Если $\rho>3/4$ и $\|p\|=\Delta$, то $\Delta$ удваивается.
  \item Если $\rho>1/4$, $\Delta$ сохраняется.
  \item Если $\rho<1/4$, $\Delta$ делится на четыре.
  \item Шаг отвергается при слишком малом $\rho$; иначе принимается.
\end{itemize}

\problemhead{B1}{Dogleg на вытянутой квадратичной функции}{У доски}
\textbf{Условие.} Для $f(x,y)=\tfrac{1}{2}(x^2+10\,y^2)$ в $x_0=(2;\ 1)$ при $B_0=\nabla^2 f$ и $\Delta=2$ найдите точку Коши $p_C$ и точку Ньютона $p_N$. Точка Ньютона лежит за шаром, поэтому dogleg идёт по ломаной из $0$ через $p_C$ в $p_N$ до пересечения со сферой. Дайте вторую компоненту $p_C$, первую компоненту шага dogleg и значение $f$ после шага.

\textbf{Решение.}
\begingroup\sloppy
\begin{enumerate}[label=\textbf{\arabic*.}, leftmargin=*]
  \item $g=Bx_0=(2;\ 10)$, $B=\operatorname{diag}(1;\ 10)$. Шаг Ньютона $p_N=-B^{-1}g=(-2;\ -1)$ имеет длину $\|p_N\|=2{,}2361>\Delta$. Безусловный минимум уже лежит за шаром, значит, dogleg остановится на границе.
  \item Квадратичная $m(-\alpha g)=f-\alpha g^\top g+\tfrac{1}{2}\alpha^2 g^\top Bg$ минимизируется при $\alpha^*=g^\top g/(g^\top Bg)=0{,}1036$; её длина $\alpha^*\|g\|=1{,}0564<\Delta$, поэтому $p_C=-\alpha^*g=(-0{,}2072;\ -1{,}0359)$. Обе координаты сжаты: крутая сильнее, чем пологая.
  \item На отрезке от $p_C$ до $p_N$ решим $\|p(\tau)\|^2=\Delta^2$ для положительного корня: $\tau=0{,}8488$. Шаг dogleg равен $p=(-1{,}7289;\ -1{,}0054)$, его длина $\|p\|=2=\Delta$, на границе.
  \item $f(x_0)=7$, $f(x_0+p)=0{,}0369$, модельная убыль $m(0)-m(p)=6{,}9631$, значит, $\rho=1>3/4$. Шаг достиг границы, поэтому $\Delta$ увеличивается; радиус следующей пробы — $4$.
\end{enumerate}
\endgroup
\begin{answerbox}
$p_C=(-0{,}2072;\ -1{,}0359)$, шаг dogleg $p=(-1{,}7289;\ -1{,}0054)$, $f=0{,}0369$.
\end{answerbox}

\problemhead{B2}{Доверительная область сжимается, когда модель лжёт}{Самостоятельно}
\textbf{Условие.} Для $f(x)=-x+\tfrac{1}{4}x^4$ в $x_0=0{,}5$ при $B_0=f''(x_0)$ начните с $\Delta_0=2$ и запустите метод доверительной области до первого принятого шага. Правило: при $\rho<0{,}25$ шаг отвергается и $\Delta\leftarrow\Delta/2$; иначе принимается. Укажите $\Delta$ принятого шага, следующую точку и число отказов.

\textbf{Решение.}
\begingroup\sloppy
\begin{enumerate}[label=\textbf{\arabic*.}, leftmargin=*]
  \item В точке $x_0=0{,}5$: $g=-0{,}875$, $B=0{,}75$. Безусловный минимум модели равен $-g/B=1{,}1667$: он лежит далеко за первым шаром радиуса $2$.
  \item Первая проба имеет длину $2$ и идёт в сторону $-\operatorname{sign}(g)$. Квартичное слагаемое растёт настолько быстро, что либо $f(x_0+p)>f(x_0)$, либо $\rho<0{,}25$; $\Delta$ делится на $2$, и проба повторяется.
  \item Принятое $\Delta=0{,}5$ приходит после $2$ отказов. Шаг равен $0{,}5$, $x_1=1$; $f$ падает с $-0{,}4844$ до $-0{,}75$.
\end{enumerate}
\endgroup
\begin{answerbox}
$\Delta=0{,}5$, $x_1=1$; отказов — $2$.
\end{answerbox}

% =====================================================================
\section{BFGS с ограниченной памятью}\label{sec:lbfgs}

Плотный BFGS хранит матрицу $n\times n$. \textbf{L-BFGS} хранит последние $m$ пар $(s_i,y_i)$ и умножает $G_k$ на $g_k$, не строя её: с масштабированием $\gamma_k=s_{k-1}^\top y_{k-1}/(y_{k-1}^\top y_{k-1})$ и двухпроходной рекурсией одно направление стоит около $4mn$ умножений со сложением.

\begin{itemize}
  \item Положите $q=g_k$. Для пар от новой к старой: $a_i=s_i^\top q/(y_i^\top s_i)$, $q\leftarrow q-a_iy_i$.
  \item Масштабируйте: $r=\gamma_kq$.
  \item Для пар от старой к новой: $b_i=y_i^\top r/(y_i^\top s_i)$, $r\leftarrow r+s_i(a_i-b_i)$.
  \item Верните $p_k=-r$: около $4mn$ операций и $2mn$ хранимых чисел.
\end{itemize}

\begin{theorembox}{Утверждение.}
Результат равен $-Gg$ для матрицы $G$, полученной из $\gamma_kI$ последовательными $m$ обновлениями BFGS по хранимым парам. При $m\ge n$ L-BFGS и плотный BFGS совпадают с точностью машинного эпсилон.
\end{theorembox}

\problemhead{C1}{BFGS вручную}{У доски}
\textbf{Условие.} Пусть $f(x)=\tfrac{1}{2}x^\top Qx$, $Q=\operatorname{diag}\begin{pmatrix} 1 & 0 \\ 0 & 4 \end{pmatrix}$, $x_0=(2;\ 2)$, $G_0=I$, линейный поиск точный. Сделайте два шага: после первого составьте $s$, $y$ и обновите $G$ по BFGS в обратной форме, проверьте секущее условие, затем сделайте второй шаг. Укажите диагональ $G_1$, $\alpha_1$ и $x_2$.

\textbf{Решение.}
\begingroup\sloppy
\begin{enumerate}[label=\textbf{\arabic*.}, leftmargin=*]
  \item $g_0=Qx_0=(2;\ 8)$, $p_0=-g_0$, $\alpha_0=g_0^\top g_0/(g_0^\top Qg_0)=0{,}2615$. Отсюда $s=\alpha_0p_0=(-0{,}5231;\ -2{,}0923)$, $x_1=(1{,}4769;\ -0{,}0923)$, $g_1=Qx_1=(1{,}4769;\ -0{,}3692)$ и $y=g_1-g_0=(-0{,}5231;\ -8{,}3692)$. На квадратичной функции $y=Qs$ точно.
  \item $y^\top s=17{,}7846>0$, $\rho=1/(y^\top s)=0{,}0562$. При $V=I-\rho sy^\top$
\handoutdisplay{
G_1=VG_0V^\top+\rho ss^\top=\begin{pmatrix} 1{,}0454 & -0{,}0028 \\ -0{,}0028 & 0{,}2502 \end{pmatrix}.
}

  \item Проверки: $G_1y=(-0{,}5231;\ -2{,}0923)=s$, $G_1^\top=G_1$, собственные значения $(0{,}2502;\ 1{,}0455)$, оба положительны. Обратная $B_1=\begin{pmatrix} 0{,}9566 & 0{,}0109 \\ 0{,}0109 & 3{,}9973 \end{pmatrix}$ отличается от $Q$ в направлении, ортогональном $s$: одна секущая пара не фиксирует всю матрицу.
  \item $p_1=-G_1g_1=(-1{,}5451;\ 0{,}0966)$, $g_1^\top p_1=-2{,}3176<0$, точный шаг $\alpha_1=0{,}9559$: $x_2=(0;\ 0)=x^*=0$. При точном линейном поиске BFGS минимизирует $n$-мерную положительно определённую квадратичную функцию не более чем за $n$ шагов.
\end{enumerate}
\endgroup
\begin{answerbox}
$G_1=\begin{pmatrix} 1{,}0454 & -0{,}0028 \\ -0{,}0028 & 0{,}2502 \end{pmatrix}$, $\alpha_1=0{,}9559$, $x_2=(0;\ 0)$.
\end{answerbox}

\problemhead{C2}{L-BFGS: двухпроходная рекурсия вручную}{Самостоятельно}
\textbf{Условие.} В задаче из C1 храните последние $m=2$ пар. Сделайте два шага L-BFGS из $x_0=(2;\ 2)$ с точным линейным поиском, применяя двухпроходную рекурсию с масштабированием $\gamma_k=s_{k-1}^\top y_{k-1}/y_{k-1}^\top y_{k-1}$. Укажите $\gamma_1$, первую компоненту направления $p_1=-H_1g_1$, вычисленного рекурсией, и $x_2$.

\textbf{Решение.}
\begingroup\sloppy
\begin{enumerate}[label=\textbf{\arabic*.}, leftmargin=*]
  \item Нулевой шаг — $-g_0$ с точным $\alpha$: он совпадает с C1. После него $\gamma_1=s_0^\top y_0/(y_0^\top y_0)=0{,}2529$.
  \item Двухпроходный проход с парами $[(s_0,y_0)]$ и ядром $\gamma_1 I$ даёт $p_1=(-0{,}3908;\ 0{,}0244)$. Точный шаг $\alpha_1=3{,}7794$, значит, $x_2=(0;\ 0)$. На 2-мерной квадратичной функции L-BFGS с $m=n$ совпадает с плотным BFGS до знака.
  \item Двухпроходная рекурсия не строит $G$: она вычисляет $Gg$ за $O(mn)$ операций. При $m\ll n$ она хранит усечённую по последним направлениям версию обратного гессиана.
\end{enumerate}
\endgroup
\begin{answerbox}
$\gamma_1=0{,}2529$, первая компонента $p_1=-0{,}3908$, первая компонента $x_2=0$.
\end{answerbox}

\begin{trapbox}
Выбирайте $m$ по задаче: слишком малое перезапускает метод каждые $m$ шагов и теряет итерации; слишком большое тратит работу без выигрыша. Обычный диапазон — от пяти до двадцати, а двухпроходной рекурсии полезно, чтобы $s^\top y$ новейшей пары оставалось положительным.
\end{trapbox}

\end{document}
